% v223: synchronized and condensed 49-item Japanese response, 77 marked-PDF anchors.
% Full Japanese translation of response_round1.tex (not a journal-submission file).
% Text is split into included TeX parts for maintainability.
% Manuscript v221: 27-page marked manuscript, 49 responses, 77 verified anchors.
% Source freshness is enforced by scripts/sync_and_check.py.
\documentclass[10pt,a4paper]{ltjsarticle}
\usepackage{luatexja-fontspec}
% Use the TeX Live default Japanese fonts (portable across CI runners).
\usepackage{amsmath,amssymb,booktabs,longtable,array}
\usepackage{stmaryrd}
\newcommand{\NT}[1]{\llbracket #1\rrbracket}
\usepackage[margin=20mm]{geometry}
\usepackage[hidelinks,unicode]{hyperref}
\usepackage{bookmark}
\usepackage{xcolor}
\setlength{\parindent}{1em}
\setlength{\parskip}{3pt}
\setlength{\emergencystretch}{3em}
\newcommand{\Comment}[1]{\par\medskip\noindent\textbf{#1　査読コメント（引用部分の全訳）：}\ }
\newcommand{\Reply}{\par\noindent\textbf{著者回答：}\ }
\newcommand{\Location}[1]{\par\noindent\textbf{修正箇所：} #1\par}
\newcommand{\ResHead}[1]{\subsubsection*{#1}}
\newcommand{\Chcf}{\mathcal C_h^{\mathrm{cf}}}
\newcommand{\Chlin}{\mathcal C_h^{\mathrm{lin}}}
\newcommand{\Bh}{\mathsf{Rec}_h}
\newcommand{\Ah}{\mathsf{Learn}_h}
\newcommand{\RSh}{\mathsf{RS}_h}
\newcommand{\RS}{\mathsf{RS}}
\title{査読者への回答 --- TCS-D-26-00494\\\large 日本語全訳（著者確認用）}
\author{Takayuki Kuriyama}
\date{}
\begin{document}
\maketitle
\noindent\textbf{論文：}Distributional Learning of Context-Free Languages under Fixed Finite-Monoid Typing

\noindent\textbf{注記：}本書は英語版 \texttt{response/response\_round1.tex} の非提出用日本語訳である。引用された査読意見も省略部分の記号を保持して訳した。節・命題番号は改訂稿を指し、原稿の旧番号は「投稿時」と明記した。ページ・余白行番号表は英語版から自動的に同期しており、最新版の変更表示付きPDFで49項目・77参照の位置を再検証した。英語本文に変更があった場合は同期検査により未更新を検出する。


% --- BEGIN japanese/parts/00_overview.tex ---
% v223: synchronized and condensed English-response overview; 49-item locator table is generated separately.

編集者および二名の査読者の詳細かつ建設的な報告に感謝する。指摘を受け、学習アルゴリズムを修正し、未使用の型付け成分を削除し、終端規則に関する完全性証明の不備を補修し、原稿の構成を大幅に改めた。

本回答書の節・定理類の番号は特記しない限り改訂稿のものとし，投稿時の番号は「投稿時」または「旧」と明記する。「第7節」「第8--10節」など査読者の見出しに沿った小見出しは投稿時原稿を指す。編集者の指摘をE.1，査読者1の指摘をR1.1--R1.37，査読者2の指摘をR2.1--R2.11とした。引用は体裁を整え，省略は省略記号で示す。下の参照表は修正表示PDFで確認されたページと余白行番号を示す。

\section*{改訂点の概観}
\noindent\textbf{中心的な修正。}
第一に，集合依存再構成作用素 $\Bh$ と保守的逐次学習器 $\Ah$ を分け，Gold学習器の構文的安定化を証明した。第二に，不要な外側文脈型注釈を削除し，仮説文法を観測部分語 $x$ ごとの $\NT{x}$ に単純化した。バッチ構成はYoshinakaの部分語文法に沿い，単項規則で窓比較を $h$-型の一致に置き換える。第三に，型付き終端規則の証拠語を明示し，必要に応じて語彙規則(L)の前に単項規則(U)を適用して完全性証明の欠落を修正した。第四に，固定 $h$ の代入可能性を非空内部因子について述べ，$\lambda$ は開始記号で別に扱う。

\noindent\textbf{先行研究と貢献。}
序論では2014年の著者プレプリントを歴史的背景とし，Yoshinaka (2015) の多ソート枠組みが固定 $h$ の定性的学習可能性も扱うことを認める。その定性的な存在，Pinによる局所自明半群の判定条件，固定窓シグネチャは新規性としない。本稿の貢献は符号長で測る特徴データ上界，型付け変更に関する閉包性，分離結果にある。Kanazawaの範疇文法の学習およびEilenbergによる正規言語の有限モノイド特徴づけは別の先行結果として帰属させた（R1.3，R2.1，R2.5）。

\noindent\textbf{主張範囲の限定。}
提示順序独立性と所定の定量的条件下での多項式特徴データ上界は集合依存作用素 $\Bh$ に限る。逐次学習器 $\Ah$ はGold収束と多項式時間更新を満たすが，逐次的な多項式データ保証は主張しない（第2.1節，R1.1）。投稿時第7節の表題はこの区別に沿って撤回した（R2.3）。外側文脈の型が必要という旧主張も撤回し，独立の文脈輸送規則を削除した（R2.2）。

\noindent\textbf{新しい定量的結果。}
命題6.3はサイズと通常厚みが $O(n)$ の有限Clark--Eyraud代入可能言語族 $L_n$ を与える。自明型付けの $\Bh$ には多項式特徴データがあるが，固定二元半束型付け $h_{\mathrm U}$ では任意の特徴標本に長さ $3\cdot2^n-2$ の語が必要となる。命題6.3(iv)は区別すべき二標的の双方に特徴標本をもつ任意の集合依存作用素へ下界を拡張する。系6.4は任意の縮約済みSSBNF表示 $G$ に対して $\tau_{h_{\mathrm U}}^{\mathrm{typ}}(G)\ge3\cdot2^n-2$ を与える。補題7.6は議論を完結させる初等的な半群論の事実として $U_1\notin\mathbf{LI}\vee\mathbf{Grp}$ を示す。したがって系7.5の代数的仮定を一様な特徴データ上界から単純には外せない。ただし任意の型付けについての必要性，任意のGold学習器の実行時間・収束時刻下界は主張しない（R2.1）。

\noindent\textbf{再構成と明確化。}
定義は第2節にまとめ，再構成作用素を完全性証明用の標的精密化より先に置いた。第6--8節は定量的議論を整理し，査読者2の線形定理への問いに答えた（R2.3）。旧三つの境界関連の節は第9節に統合し，命題9.7を加えた。正規化証明は付録A--Bへ移し，付録Cでは可逆な決定的解析器を用いて命題6.3(i)を証明した。投稿時付録Aの閉包性は命題3.2に簡約した（R1.36）。正規化後の線形文法 $G_0$ に対する定理8.3の符号長評価 $\|\mathsf{Wit}(\widetilde G_0)\|=O(|G_0|^3|M|^4)$ は明示したが，正規化前の文法に同じ次数を主張しない。非消去逆像（第3節），分解 $\mathsf{S}^n=\mathsf{S}E(\mathsf{S})\mathsf{S}$（命題3.5），補題7.1の重要節点の幅の評価も説明した。命題9.1と $\rho\ge2$ の鋭さの変更は R1.33--R1.35 に記載した。

\subsection*{投稿時から改訂後への主要番号対照}
原稿の構成を大幅に変えたため、報告書で頻繁に言及された項目の対応を示す。
\begin{center}\small
\begin{tabular}{@{}p{0.46\linewidth}p{0.46\linewidth}@{}}
\toprule
投稿時の原稿 & 改訂原稿\\\midrule
定義2.3（$\sim_h$）& 定義2.1\\
定義3.2（Clark--Eyraud）& 第3.1節の番号なしの特殊例\\
命題3.5 & 命題3.3\\
定理3.7 & 定理3.4\\
例4.3 & 削除\\
旧規則(2),(3) & 商構成と単項規則(U)\\
補題5.2(i) & 証拠集合(2)、定理5.3\\
定理5.7 & 定理5.4\\
系5.8 & 系5.5\\
命題7.2 & 命題8.1\\
定理7.9 & 定理8.3\\
定理8.6 & 定理9.2（証明簡略化）\\
系8.7 & 削除\\
命題10.3--10.5 & 命題9.1と命題3.3直後の判定条件\\
補題9.11 & 補題9.6\\
付録A（閉包性）& 命題3.2（規約変更に対応）\\
$\mathrm{CCL}_{p},\mathrm{CCL}$ & $\mathrm{CTR}_{\rho},\mathrm{CTR}$\\
カウンター記号 $d,u,;$ & $d\mapsto\uparrow$（増加），$u\mapsto\downarrow$（減少），$;\mapsto\#$（リセット）\\
$L_Y$（古典的言語）& $L_{\mathrm{Luk}}$（Łukasiewicz言語）\\
命題9.2, 9.6, 9.9・補題9.8 & 第9.3節の短い記述・引用\\
\bottomrule
\end{tabular}
\end{center}

% --- END japanese/parts/00_overview.tex ---


% --- BEGIN japanese/locator_table_JP.tex ---
\subsection*{修正表示PDFに対する参照表}
以下の49項目・77参照は、英語版査読回答の検証済み参照表から自動的に転記される。
\begingroup\small
\setlength{\LTpre}{4pt}\setlength{\LTpost}{6pt}
\begin{longtable}{@{}p{0.09\linewidth}p{0.83\linewidth}@{}}
\toprule
項目 & 改訂原稿での位置\\
\midrule\endfirsthead
\toprule
項目 & 改訂原稿での位置\\
\midrule\endhead
E.1 & 定義 2.1 （Marked PDF p.~4，余白行~169）; 第4--5節 （Marked PDF p.~8，余白行~324）\\
R1.1 & 第2.1節 （Marked PDF p.~4，余白行~135）; 第4.2節 （Marked PDF p.~9，余白行~361）; 第5.2節 （Marked PDF p.~11，余白行~437）\\
R1.2 & 序論 （Marked PDF p.~1，余白行~23）; 付録 A--C （Marked PDF p.~23，余白行~885）\\
R1.3 & 序論と関連研究 （Marked PDF p.~1，余白行~23）\\
R1.4 & 序論 （Marked PDF p.~1，余白行~23）; 定義 2.1 （Marked PDF p.~4，余白行~169）\\
R1.5 & 第1.2節 （Marked PDF p.~2，余白行~80）\\
R1.6 & 第9.3--9.4節 （Marked PDF p.~20，余白行~810）\\
R1.7 & 第4.2節 （Marked PDF p.~9，余白行~361）\\
R1.8 & 第2節 （Marked PDF p.~3，余白行~110）; 第4.1節 （Marked PDF p.~8，余白行~325）\\
R1.9 & 第5.1節 （Marked PDF p.~10，余白行~376）\\
R1.10 & 序論 （Marked PDF p.~1，余白行~23）; 第3節 （Marked PDF p.~5，余白行~188）\\
R1.11 & 補題 2.2 （Marked PDF p.~4，余白行~169）; 第3.1節 （Marked PDF p.~6，余白行~244）\\
R1.12 & 第3.1節 （Marked PDF p.~6，余白行~244）; 命題 3.3 （Marked PDF p.~6，余白行~255）; 命題 3.5 （Marked PDF p.~8，余白行~303）\\
R1.13 & 第2節 （Marked PDF p.~3，余白行~110）; 定理 3.4 （Marked PDF p.~7，余白行~279）\\
R1.14 & 定理 3.4 （Marked PDF p.~7，余白行~279）\\
R1.15 & 第2節 （Marked PDF p.~3，余白行~110）; 第5.1節 （Marked PDF p.~10，余白行~376）\\
R1.16 & 第2節 （Marked PDF p.~3，余白行~110）\\
R1.17 & 第5.1節 （Marked PDF p.~10，余白行~376）\\
R1.18 & 第4節 （Marked PDF p.~8，余白行~324）; 第5.1--5.2節 （Marked PDF p.~10，余白行~376）\\
R1.19 & 命題 5.2 （Marked PDF p.~10，余白行~412）\\
R1.20 & 第5.1節 （Marked PDF p.~10，余白行~376）\\
R1.21 & 第5.1節 （Marked PDF p.~10，余白行~376）\\
R1.22 & 命題 5.2 （Marked PDF p.~10，余白行~412）\\
R1.23 & 命題 5.2 （Marked PDF p.~10，余白行~412）\\
R1.24 & 第4.2節 （Marked PDF p.~9，余白行~361）; 第5.2節 （Marked PDF p.~11，余白行~437）\\
R1.25 & 第4.2節 （Marked PDF p.~9，余白行~361）; 第5.2節 （Marked PDF p.~11，余白行~437）\\
R1.26 & 第4.1節 （Marked PDF p.~8，余白行~325）\\
R1.27 & 定理 4.2 （Marked PDF p.~9，余白行~348）\\
R1.28 & 定理 5.3 （Marked PDF p.~11，余白行~438）\\
R1.29 & 定理 5.4 （Marked PDF p.~12，余白行~459）\\
R1.30 & 定理 5.6 （Marked PDF p.~12，余白行~480）\\
R1.31 & 第8節 （Marked PDF p.~17，余白行~662）; 付録 B （Marked PDF p.~24，余白行~950）\\
R1.32 & 第8節 （Marked PDF p.~17，余白行~662）; 付録 A--B （Marked PDF p.~23，余白行~885）\\
R1.33 & 定理 9.2 および 第9.2節 （Marked PDF p.~20，余白行~803）\\
R1.34 & 第9節 （Marked PDF p.~19，余白行~743）\\
R1.35 & 命題 9.1 （Marked PDF p.~19，余白行~768）; 第9.2--9.3節 （Marked PDF p.~20，余白行~795）\\
R1.36 & 補題 9.6 （Marked PDF p.~21，余白行~829）\\
R1.37 & 命題 9.7 （Marked PDF p.~21，余白行~844）\\
R2.1 & 序論 （Marked PDF p.~1，余白行~23）; 第3節 （Marked PDF p.~5，余白行~188）; 命題 6.3 （Marked PDF p.~14，余白行~540）; 補題 7.6 （Marked PDF p.~16，余白行~648）; 第8--9.2節 （Marked PDF p.~17，余白行~662）\\
R2.2 & 第4.1節 （Marked PDF p.~8，余白行~325）; 第5.1--5.2節 （Marked PDF p.~10，余白行~376）\\
R2.3 & 定理 8.3 （Marked PDF p.~17，余白行~683）; 命題 8.4 （Marked PDF p.~18，余白行~702）\\
R2.4 & 第4--5節 （Marked PDF p.~8，余白行~324）; 付録 A--B （Marked PDF p.~23，余白行~885）\\
R2.5 & 命題 3.1 （Marked PDF p.~5，余白行~194）\\
R2.6 & 第5.1節 （Marked PDF p.~10，余白行~376）\\
R2.7 & 第5.1節 （Marked PDF p.~10，余白行~376）\\
R2.8 & 第5.1節 （Marked PDF p.~10，余白行~376）; 定理 5.3 （Marked PDF p.~11，余白行~438）\\
R2.9 & 定理 5.3--5.4 （Marked PDF p.~11，余白行~438）\\
R2.10 & 系 5.5 （Marked PDF p.~12，余白行~465）\\
R2.11 & 命題 8.1 （Marked PDF p.~17，余白行~666）; 補題 8.2 （Marked PDF p.~17，余白行~673）; 付録 B （Marked PDF p.~24，余白行~950）\\
\bottomrule
\end{longtable}
\endgroup

% --- END japanese/locator_table_JP.tex ---


% --- BEGIN japanese/parts/01_editor_r1_01_14.tex ---

\section*{修正表示原稿の参照表について}
上の英語版から自動同期される表は、各意見に対応する改訂後の節・命題と、実際の修正表示PDFで検証したページ・余白行番号を示す。節への参照は複数の改訂段落を含みうるが、定理類の項目名は対象となる記述を特定する。

\section*{編集者への回答}
\Comment{E.1}
「査読者はともに本論文の科学的貢献に価値があると認めている。しかし両者とも提示方法に多くの問題を指摘しており、その多くは単なる文章の修正にとどまらず、論文の科学的部分にも及んでいる。学習に関する部分では特にそうである。……とくに正しさに関する問題は慎重に検討すべきである。」

\Reply
同意する。投稿時の学習器を保守的な Gold 学習器に置き換え、終端規則に関する欠落を証拠集合と完全性証明の双方で補修し、使用されていない外側文脈の注釈を削除した。固定 $h$ の代入可能性は非空の内部因子について定義し、$\lambda$ は別途扱う。第4節でバッチ作用素と逐次学習器を定義した後、第5節で完全性のための標的側の生成語型付けを導入する。多項式特徴データは集合依存バッチ作用素に関する主張であり、逐次学習器については多項式時間更新を伴う Gold の安定化を示す。
\Location{定義2.1、第4.2節、第5.1節と定理5.3、系5.5、第6--8節。}

\section*{Reviewer 1}
\subsection*{総評と学習アルゴリズム}

\Comment{R1.1}
「本投稿の主結果は、Clark and Eyraud (2007) の代入可能文脈自由言語（SCFL）の学習、および Yoshinaka (2008) の $(k,\ell)$-SCFL の学習を洗練された形で一般化したものである。……唯一の重大な誤りは学習アルゴリズムの定義にあるが、先行研究に従えば容易に修正できる。」

\Reply
中心的な正しさの問題を指摘されたことに感謝する。正例ごとにバッチ文法を再構成する旧方式では構文的安定化が保証されないため，集合依存の再構成作用素 $\Bh$（第4.1節）と保守的な逐次学習器 $\Ah$（第4.2節）を分離した。後者は現在の仮説が生成しない正例を受け取った場合にのみ再構成し，系5.5が最終的な構文的安定化を証明する。Eyraud--Heinz--Yoshinaka (2016, Appendix) の遅延再構成の現象も引用した。したがって逐次学習器に主張するのは Gold 収束と多項式時間更新であり，逐次的な多項式時間・データ同定（同 Definition 2.5）は主張しない。多項式特徴データの保証は集合依存の $\Bh$ に対するもので（同 Definition 2.8），明示した定量的範囲に限る。第2.1節でこの区別を説明した。
\Location{第2.1節、第4.2節、系5.5、定理7.2、定理8.3およびその後の補足。}

\Comment{R1.2}
「一方、提示にはかなり改善の余地がある。技術用語の定義の欠如、不必要に長い証明、メタ変数の使い方の不統一、節を過度に細分化している点を直すことで、読者の負担を大幅に減らせるはずであり、そのような大改訂を勧める。……この内容は現在の半分以下の長さで提示できるのではないか。」

\Reply
提示を大幅に短縮した。投稿時原稿は40ページであり，付録・参考文献を含む改訂稿はClean版・修正表示版ともに27ページである。説明の反復を除き，証明義務を残して序論・結論・線形正規化を短縮した。番号付き定理類は半数以下，本文の節数は11から10，小節数は31から16となった。第2節に定義と正例学習モデルをまとめ，旧計算量節を第5節に統合した。第6--8節では定量的議論を整理し，旧三つの境界関連の節は短縮した第9節に統合した。

未定義の用語とメタ変数の不統一を修正し（R1.8，R1.11），SSBNFの説明は本文に置いた。両査読者の提案に従い，型付き生成語と正規化の証明を付録A，線形spineの正規化を付録B，命題6.3(i)の証明を付録Cに配置した。
\Location{第1--2節、第4--9節、付録A--C。}

\subsection*{投稿時第1節}
\Comment{R1.3}
「研究分野内での位置づけと対象言語族の重要性を述べれば、論文の魅力が増す。一方、現状は技術的な内容に重きを置きすぎているので、より簡潔にすべきである。」

\Reply
同意する。第1.2節では固定 $h$ の定性的な学習構成と Yoshinaka の一般的枠組みとの関係を簡潔に述べ，貢献を定量的保証と構造・表現力上の位置づけの二つに整理した。固定準同型は有限の型付けバイアスであり，学習を主張する固定 $h$ のクラスと，表現力の比較に使う $\RS$ を区別した。2014年のプレプリントと Yoshinaka (2015) との関係，固定有限モノイド型付けに固有の定量的・代数的結果も明示した（R2.1を参照）。
\Location{第1節。}

\Comment{R1.4}
「$\mathcal D_L$ と $\mathcal C_h^{\mathrm{cf}}$ は標準的な記号ではないので定義が必要である。とくに『一般の文脈自由クラス $\mathcal C_h^{\mathrm{cf}}$』という言い方は誤解を招きやすい。」

\Reply
同意する。$\mathcal D_L(x)$ は序論での初出時と予備事項で定義した。$\mathcal C_h^{\mathrm{cf}}$ も序論の初出時に $\sim_h$-代入可能な文脈自由言語のクラスとして定義する。$h$ を固定した制約のないクラス $\RSh$ は、合併 $\RS=\bigcup_h\RSh$ より先に導入した。第3節でこれらのクラスと $\Chlin$ を形式的に扱い、「一般の文脈自由クラス」という表現は削除した。
\Location{序論、第2.2節および第3節冒頭。}

\Comment{R1.5}
「貢献が(1)--(6)と列挙されているが、細かな技術的事項まで独立の貢献として数える必要はない。また、この段階では専門用語の意味が未だ不明なものが多く、読者はついていけない可能性がある。」

\Reply
同意する。第1.2節では定性的学習構成と Yoshinaka の枠組みを番号なしの短い段落で紹介し，番号付き貢献を定量的保証と構造・表現力上の位置づけの二つに絞った。定量的貢献では特徴データの多項式上界と指数下界を概説し，「型付き厚み」と「開始記号分離二分標準形」に言及するが，それぞれの形式的定義は第6節と第5.1節で与える。構造上の貢献では固定窓の判定条件を古典的結果として示し，「局所自明」は第2.2節で定義した。第1.1節の $L=\{a,aa\}$ の例は，自明型付けでは代入可能性が破れる一方，偶奇型付けで因子を区別できることを示す。これは定義の例であり，固定窓クラスからの分離例ではない。
\Location{第1.1節（固定 $h$ 代入可能性の定義直後の最小例）、第1.2節。}

\Comment{R1.6}
「『Yoshinaka の古典的言語』は、文献では『Łukasiewicz 言語』と呼ぶことが多い。」

\Reply
同意する。「Łukasiewicz 言語」を用い、Autebert--Berstel--Boasson (1997) のハンドブック文献を加え、1種類の括弧からなる Dyck 言語との関係を簡潔に述べた。この言語がすべての固定 $(k,\ell)$ クラスの外にあるという既知の観察は Yoshinaka (2008, Section 3) に帰属させ、有限型付けの合併 $\RS$ の外にあるという本論文の強い結論と区別した。
\Location{第9.3節。}

\subsection*{投稿時第2節}
\Comment{R1.7}
「『この論文で考える学習器は、正例の提示順序によらず同じ出力を生じる』という主張は疑わしい。学習アルゴリズムと、その部分手続きである文法構成アルゴリズムが混同されているように見える。実際、第5節の学習アルゴリズムは収束せず、Gold の基準を満たさない。Clark と Eyraud に従って修正しても、例えば $a^ncb^n$ の学習では、特徴標本を最初に受け取る場合と、最初の正例が非常に長い場合とで、最終出力は異なるだろう。」

\Reply
同意する。投稿時の順序独立性の主張を撤回し，集合依存のバッチ作用素と保守的逐次学習器を区別した。第4.2節では任意の無限標的について，十分長い正例をテキストの先頭に追加すると，安定後にも余分な観測非終端記号が残りうることを示し，提示順序によって最終文法が異なると説明する。Clark (2013, Section 4, Example 1) にも同様の例がある。
\Location{第2.1節、第4.2節、系5.5。}

\Comment{R1.8}
「この節では、（線形）文脈自由文法を含め、必要な記法と用語を十分に定義すべきである。『type』『reconstruction』『realize』『reachable』『observation』『control』『local configuration』などが脚注で場当たり的に定義されている。数学的定義のない直感的な言葉に過ぎないなら、名前を付けないほうがよい場合もある。過剰な用語の氾濫は読者の負担になり、理解を妨げる。」

\Reply
同意する。第2節に線形CFG，到達可能・生成可能・縮約済み文法，正例学習モデル，再構成，特徴標本，分布，固定 $h$ のクラスをまとめた。証明内部の不要な術語を削除し，部分語に対応する仮説非終端記号 $\NT{x}$ は第4.1節で初出時に定義した。文字数 $|w|_a$，SGL と $(k,\ell)$-SGL の略称，$\mathsf{Acc}\subseteq M$ も初出時に説明し，未使用の $[x]_{\equiv_L}$ を削除した。列挙された語のうち，\emph{reachable} と生成語の \emph{type} は各定義箇所で，\emph{reconstruction} はバッチ作用素として定義する。\emph{observation} は独立の術語にせず，\emph{realize} と \emph{local configuration} は削除した。\emph{control} は制御集合などの標準的な用法に限る。
\Location{第2節、第4.1節、第5.1節、付録A.2。}

\Comment{R1.9}
「『finite typed reconstruction basis』という語はここでは定義されているが、その四つの単語が個別には未定義なのに全体が不可分の概念なら、もっと短い名称のほうがよい。」

\Reply
同意する。名称を変えるのではなく、証明内部のこの対象と『基底』自体を削除した。新しい完全性証明は、生成語型付きの標的文法と厳密な再構成に必要な有限証拠集合 $\mathsf{Wit}(\widetilde G)$ のみを用いる。
\Location{第5.1節。}

\subsection*{投稿時第3節}
\Comment{R1.10}
「論文は $h$ が固定されることを、やや過度かつ不自然な形で強調している。それに合わせるなら $\mathsf{RS}_h$ を $\mathsf{RS}$ より先に定義すべきだ。$\mathsf{RS}$ が正例から学習可能だとは示されておらず、実際それは学習可能でない。学習可能なのは $\mathsf{RS}_h$ である。」

\Reply
同意する。序論で $\RSh$ を $\RS:=\bigcup_h\RSh$ より先に定義し、Gold 同定定理が扱うのは、固定 $h$ に対する $\Chcf=\RSh\cap\mathrm{CFL}$ の非空言語であって、$\RSh$ 全体でも変動する $h$ に関する合併 $\RS$ でもないことを明示した。後者は表現力の比較にのみ使う。また、$\RS$ と $\RS\cap\mathrm{CFL}$ は非空アルファベット上で正例から同定可能ではない。すべての有限言語と $\Sigma^*$ を含むためである。
\Location{要旨、序論、第2.2節および第3節冒頭。}

\Comment{R1.11}
「定義3.2は注意を要する。Clark and Eyraud (2007) の定義の述べ方はいくらか不注意だが、同論文第6.1節で $a^+$ が代入可能であると述べ、文法構成アルゴリズムで空語を除外していることを考えると、Yoshinaka (2008) の $(0,0)$-代入可能性を意図したと解するのが妥当だろう。……」

\Reply
査読者による Clark--Eyraud (2007) の解釈に同意する。非空内部因子に制限すれば意図された場合は Yoshinaka の $(0,0)$ と一致する（第3.1節）。Clark (2013) を引用し，Clark--Eyraud の Definition 5 が一般の文字列を量化する点と，改訂稿の非空因子の規約を区別した。固定 $h$ の代入可能性は非空の内部因子に課し，$\lambda$ は開始記号で扱う。新たな補題2.2において，元のモノイドの基礎半群に別個の単位元を付加すれば，空語を元の単位元型の非空語と区別でき，空語も量化する条件が元の代入可能性と同値になることを証明した。再構成と開始記号分離標準形は変更しない。例えば $L=a^+$ では $(a,\lambda)$ は $\mathcal D_L(\lambda)$ と $\mathcal D_L(a)$ の双方に属するが，$(\lambda,\lambda)$ は後者だけに属するので，空因子の比較は $a^+$ を排除する。中心印付き言語の例でも空因子の比較は用いない。$(k,\ell)$ の定義では外側の文脈を $u_1,u_2,v_1,v_2$，固定窓を $p,q$，可変内部を $r,s$ とした。
\Location{定義2.1、補題2.2、第3.1節の Clark--Eyraud の特殊例および固定窓の定義。}

\Comment{R1.12}
「命題3.5：$\sqcup$ と $\mathrm{pref}_k$ が定義されていない。……証明が非常に長い。$T$ の積を最初に明示すれば、$\tau$ がモノイド準同型であることは一行で済む。$h_{k,\ell}$ や $M_{k,\ell}$ を使った言い換えは不必要に煩雑である。」

\Reply
固定窓の議論が冗長だったことに同意する。未定義の $\sqcup$ を削除し，接頭辞・接尾辞と短語の規約を先に定義した。命題3.3は有限像の積を直接与え，準同型性を簡潔に確認する。ただし $h_{k,\ell}$ と $M_{k,\ell}$ は後続の命題3.5，補題7.1，定理7.2，補題7.4で再利用するため保持した。短語は個別の型を持つ。固定窓の反例判定条件は命題3.3の後にまとめ，古典的な局所自明性の特徴づけ，窓・群の分解，厚みに基づく上界は命題3.5，補題7.4，系7.5として区別した。
\Location{命題3.3と直後の固定窓反例条件、命題3.5、補題7.4、系7.5。}

\Comment{R1.13}
「$h$ がデータとして与えられるというより、早い段階で『効率的に計算可能な $h:\Sigma^*\to M$ を固定する』などと述べるほうが明確だ。（線形）文脈自由言語を定義してほしい。……『明示的有限モノイド準同型』は定義されていない。……定理の主張では技術的に不要な言葉を削除すべきである。」

\Reply
同意する。有限モノイドは積表と文字像で固定し、$h(w)$ の計算方法を明示した。線形 CFG を予備事項で定義し、主定理では「固定された効率的に計算可能な準同型」という表現を用いる。
\Location{第2--3.2節。}

\Comment{R1.14}
「『The proof for fixed $h$』という表現が理解できない。『定理3.7の証明』という意味か。」

\Reply
同意する。この紛らわしい句は削除し，改訂稿では対応する定理と証明を直接記した。
\Location{第3.2節。}

% --- END japanese/parts/01_editor_r1_01_14.tex ---


% --- BEGIN japanese/parts/02_r1_15_30.tex ---

\subsection*{投稿時第4節}
\Comment{R1.15}
「定義4.2：『productive』と『reachable』が未定義である。『この節は typed local control → finite typed basis の順に進む』という表現もくだけすぎている。」

\Reply
同意する。これらの語は予備事項で定義し，型付き文法は完全な精密化の縮約部分として定義したため，到達可能性・生成可能性の繰り返しの断りは不要となった。引用された非形式的な文も削除した。
\Location{第2節、第5.1節。}

\Comment{R1.16}
「脚注2：『productive』はそれ以前に使われている。脚注ではなく第2節で定義すべきだ。」

\Reply
同意する。到達可能・生成可能を標的文法の精密化より前の第2節で定義し、従来の説明脚注は削除した。
\Location{第2節、第5.1節。}

\Comment{R1.17}
「『到達可能かつ生成可能な型付き非終端記号の集合』と書いてあるが、$\widetilde G$ にそうでない記号があるのか。読者の助けを意図した強調かもしれないが、かえって混乱を招く。」

\Reply
査読者の指摘は正しい。$\widetilde G$ は完全な生成語型付き精密化から縮約部分を取り出したものとして定義されるので、残された記号はすべて到達可能かつ生成可能である。重複する限定を削除した。
\Location{第5.1節。}

\Comment{R1.18}
「冒頭の段落では、異なる出現を異なる型付きコピーとして扱う必要がなぜあるのか、理由が全く示されていない。『何のために、なぜ』を明確にする表現が必要だ。例4.3には『型付けが必要な理由』と付くのに、理由自体が述べられていない。……説明するなら、まず学習アルゴリズムを示し、その完全性証明で『ここでは標的文法が型付けにより精密化されていると仮定し、仮説文法が型付き標的導出を模倣できることを示す』と述べたほうが自然だろう。」

\Reply
提案された順序を採用した。第4節で集合依存のバッチ再構成作用素と保守的逐次学習器を先に定義し，第5節で完全性証明のために生成語型付きの標的文法を導入する。その役割は既定義の仮説文法が型付き標的導出を模倣できることの証明である。外側文脈型付けは削除した（R2.2）。生成語型付けを使う単項規則(U)の箇所では，型付き終端規則 $X\to a$ に対し $h(\omega(X))=h(a)$，型付き二項規則に対し $h(\omega(X))=h(\omega(Y))h(\omega(Z))$ が成り立つ。
\Location{第4節、第5.1節、定理5.3。}

\Comment{R1.19}
「補題4.4：『完全な型付き精密化 $\widetilde G^{\mathrm{full}}$ において』の直後に、$G$ に関する導出の式が現れるのは紛らわしい。不要であるか、文末に移すべきだ。」

\Reply
旧補題4.4は再構成によって削除した。必要な内容は命題5.2に吸収し、言語保存と生成語型の不変量を直接述べ、元の文法と型付き精密化の導出は構文的に区別した。
\Location{命題5.2。}

\Comment{R1.20}
「補題4.8：『$X$ が $\widetilde G$ において到達可能かつ生成可能ならば』という限定も同様に読者を混乱させる。」

\Reply
R1.17の縮約済み精密化の規約によって解決した。旧補題4.8と，到達可能性・生成可能性を重ねて仮定する記述を削除した。
\Location{第5.1節。}

\Comment{R1.21}
「冒頭の段落（$\omega,\chi$ の定義）は、$\mathsf{CS}$ の定義直前に置くとつながりがよいだろう。」
「『特徴標本を次で定義する』というのは、定義2.1の条件を満たすことが直ちに明らかでない限り避けるべきだ。」
「『realizable』をここで定義するのに、それより前から使っており、『realized』は未定義である。」
「『canonical observation word』は $\mathsf{CS}(\widetilde G)$ の定義箇所で定義できるだろう。」

\Reply
同意する。標準生成語 $\omega(X)$ は縮約済み精密化の直後，初出より前に定義し，標準文脈 $\chi(X)$ は有限証拠集合の直前に置いた。$\mathsf{Wit}(\widetilde G)$ は厳密な再構成を証明するまでは証拠集合と呼び，その後に特徴標本と位置づける。未使用の基底，$NT(\cdot)$，$\mathrm{Rule}(\cdot)$，realized/realizable，exhibited などの術語は削除した。
\Location{第5.1節。}

\Comment{R1.22}
「補題4.6は、あまりにも自明な系である。系と呼び直すか削除すべきだ。」
「補題4.7はここまでの議論を理解した読者には自明であり、証明は不要だ。」
「定理4.9が何を示すのかわからない。」
「定理4.12の第一文は削除すべきだ。第二、第三文は未定義の『exhibited』のため理解できない。」

\Reply
旧補題4.6--4.7および定理4.9、4.12は削除した。型付き精密化の議論は、ひとつの正規形定義と、言語保存および生成語型の不変量を述べる命題5.2、それに証拠集合の構成に整理した。
\Location{第5.1節。}

\Comment{R1.23}
「$\alpha$ は記号ではなく文字列である。」
「『とくに $X\to YZ$ ならば』というが、その前件と帰結に関連がない。」
「補題4.5の(iii)について『$G$ の各導出は型注釈を消去する』という文は理解できない。『型注釈を消した $G$ の導出は……を与える』という意味か。」
「$\mathrm{NT}(\widetilde G)$ は9ページの $\widetilde V$ と同じではないか。」
「定義4.11の $\mathrm{Rule}(\widetilde G)$ は9ページの $\widetilde P$ と違うのか。」

\Reply
同意する。不適切な文を削除し，$NT(\widetilde G)$ と $\mathrm{Rule}(\widetilde G)$ はそれぞれ $\widetilde V$ と $\widetilde P$ に統一した。持ち上げの主張は命題5.2にまとめ（R1.19），型付き導出から型注釈を消せば元の文法の導出になると証明中に明記した。
\Location{第5.1節。}

\subsection*{投稿時第5節}
\Comment{R1.24}
「ここで定義された学習アルゴリズムは収束しない。これは論文にとって致命的な誤りである。」

\Reply
同意する。投稿時の学習アルゴリズムは収束せず、これは中心的な正しさの誤りだった。R1.1で説明したように、アルゴリズムを保守的な逐次学習器に置き換え、系5.5で構文的な安定化を証明した。
\Location{第4.2節、系5.5。}

\Comment{R1.25}
「系5.8：この議論は証明として受け入れられない。第5.1節の定義では、学習器の仮説文法は正例を受け取るたびに成長し続ける。」

\Reply
R1.1およびR1.24を参照されたい。証明だけでなく学習器自体を変更した。その結果として系5.5の安定化を得た。
\Location{第4.2節、系5.5。}

\Comment{R1.26}
「注5.1：『規則の健全性』が何を意味するのかわからない。……一文加えるなら『文法は各非終端記号 $[x:u,v]$ から $x$ と代入可能な語の集合が導出されるよう設計する』といった説明が役立つ。定義をもたない形容詞を使うより、数学的性質を簡潔に述べるほうがよい。」
「補題5.3と5.4は定義を繰り返しているだけなので削除すべきだ。」

\Reply
定義の繰り返しである注と補題を削除した。仮説文法では観測された非空因子 $x$ ごとに一つの $\NT{x}$ を置き，出現文脈ごとの記号を用いない（R2.2）。健全性は $\NT{x}\Rightarrow^*w$ ならば $x\equiv_Lw$ かつ $h(w)=h(x)$ という不変量で直接述べる（R1.27）。有限標本の非終端記号が構文合同類全体を生成するとは主張しない。
\Location{第4節（規則(B),(U),(L),(S)と健全性の不変量）。}

\Comment{R1.27}
「補題5.2：第二文の第一の主張は自明で、第二の主張は理解できない。対応する『局所配置』とは何か。『より正確には』に続く部分との関係が不明だ。」

\Reply
旧補題5.2と未定義の「対応する局所配置」は削除した。代わりに定理4.2の証明で
\[
\NT{x}\Rightarrow^*w\quad\Longrightarrow\quad x\equiv_Lw\ \text{かつ}\ h(w)=h(x)
\]
という単純な不変量を用い、導出木の高さに関する帰納法で証明する。単項規則の場合は共通の標本文脈と固定 $h$ の代入可能性を、二分規則の場合は $\equiv_L$ がモノイド合同であることを用いる。開始規則によって仮説文法の健全性が従う。
\Location{定理4.2とその健全性の不変量。}

\Comment{R1.28}
「14ページ1行目：『$G$ が SSBNF なので』は $\widetilde G$ の誤りではないか。」

\Reply
書き直しの際に修正した。完全性証明全体で、対象の文法を明確に指定した。
\Location{定理5.3。}

\Comment{R1.29}
「定理5.7：主張を数学的に書くべきだ。『有限標本が有限型付き再構成基底を露出する』とは何を意味するのか。標本は文字列の集合だが、FTRB は組である。」

\Reply
定理5.4は、必要な仮定を定理文にすべて含めて、
$\mathsf{Wit}(\widetilde G)\subseteq K\subseteq L
\Longrightarrow L(\Bh(K))=L$
という明確な含意を記す。
\Location{定理5.4。}

\subsection*{投稿時第6節}
\Comment{R1.30}
「注6.4の立場を採るなら、文法構成が多項式時間で実行できることはほとんど自明であり、この節の内容は一段落で十分なのではないか。」

\Reply
独立していた計算量の節を第5.3節の短い説明にまとめた。定理5.6では $n_K=\sum_{w\in K}(|w|+1)$ と置き，空語も数える。因子非終端記号は $O(n_K^2)$，二項・単項規則の候補はそれぞれ $O(n_K^3)$ 個であり，明示的構成時間は $O(n_K^4)$ で抑えられる。更新時間については独立した系ではなく本文中に述べる。
\Location{第5.3節。}

% --- END japanese/parts/02_r1_15_30.tex ---


% --- BEGIN japanese/parts/03_r1_31_37.tex ---

\subsection*{投稿時第7節}
\Comment{R1.31}
「これは前の節の議論の特別な場合であり、ずっと短くできる。……説明が長すぎると、何を理解すべきか、前の節と何が違うのかという焦点がぼやける。」
「補題7.3：『strict linear derivation』が定義されていない。」

\Reply
投稿時第7節に相当する部分を改訂版第8節に短縮した。本文には定理8.3で使う単一spine正規形（命題8.1）と短い生成語の上界（補題8.2）を残し，補題8.2には直後に短い証明を置いた。付録Bでは統一的なspine鎖を使って正規化を証明した。未定義だった \emph{strict linear derivation} は削除した。
\Location{命題8.1と周辺の議論、補題8.2とその証明、付録B。}

\Comment{R1.32}
「命題7.2の証明に2ページも使う必要はない。全く同じ制限が定義されているかは知らないが、これは本質的でないことが知られているいくつかの制限の組合せであり、多くの読者は自明か民間伝承的事実として受け入れるだろう。」
「補題7.5と7.6はあまりにも自明である。」

\Reply
同意する。旧2ページの正規化証明はR1.31で説明した命題8.1（証明は付録B）と補題8.2（短い証明は本文）に整理した。旧補題7.5と7.6は削除するか各証明に吸収した。付録Bでは古典的正規化操作について Autebert--Berstel--Boasson (1997) と Hopcroft--Motwani--Ullman (2007) を引用し，二項規則のラッパー子がちょうど一つというspine構造に加え，文法サイズと構成時間の多項式上界を確認した。
\Location{命題8.1と周辺の議論、補題8.2とその証明、定理8.3とその後の補足、付録B。}

\subsection*{投稿時第8--10節}
\Comment{R1.33}
「アルファベット $\{d,u,;\}$ は適切でない。とくにセミコロンは本文中では読みにくく、$u$ は他の箇所で文字列を表すことがほとんどなので混乱の原因になる。……」
「定理8.6の証明は簡単にすべきだ。証拠語を短くするために複雑な式やギリシア文字を多用しても本質ではなく、読者の負担を増す。……」

\Reply
同意する。$d\mapsto\uparrow$（増加），$u\mapsto\downarrow$（減少），$;\mapsto\#$（リセット），$p\mapsto\rho$ とし，査読者が提案した証拠語の本質的な形を採用した。定理9.2の二語 $(\uparrow\downarrow)^k\uparrow^{\rho-1}\#^\ell$ と $(\uparrow\downarrow)^k\uparrow^\rho\#^\ell$ は空文脈と固定窓の端を共有するが，先行する $\uparrow$ で分布を区別できる。命題3.3直後の判定条件により長い議論を避けた。鋭い境界 $\rho\ge2$ は $\mathrm{CTR}_1$ が $(1,1)$-代入可能であることを脚注で説明した。
\Location{第9.2節。}

\Comment{R1.34}
「この話題のために小節を一つ丸ごと使う必要はない。系8.7が何を確立するのか不明である。……命題3.5と定理8.6には意味があるが、系8.7にまとめることで重要な貢献をかえって見えにくくしている。」
「『CCL の語のブロックとは、区切り記号 $;$ で区切られた部分語のそれぞれを指すことに注意する』というのは注意書きではなく新しい定義である。しかも定理8.6の証明はブロックという語を全く別の意味で使っている。」
「第8--10節はいずれも学習対象の言語クラスの豊かさに関するものであり、一つの節に統合し、説明と証明を簡潔にすれば論文の全体像が捉えやすくなる。」

\Reply
個別に学習可能な言語族の合併が学習可能とは限らないことに同意し，旧系8.7を削除した。$\RS$ は表現力の比較にだけ使い，学習可能性は固定 $h$ の文脈自由部分クラスにのみ主張する。旧命題3.5と定理8.6に対応する命題3.3と定理9.2は別の主張として残す。三つの境界関連の節は第9節に統合した。曖昧な \emph{block} の術語を因子と高さの記述に置き換え，$\mathrm{CTR}_\rho$ の型付け $h_\rho$ を第9.2節のDFAの遷移モノイド準同型として定めた。Yoshinaka (2008, Proposition 1) の正規言語の分離結果を明記した。投稿時第10節の $\Delta^*=(\{a^nb^n:n\ge0\})^*$ は命題9.1に短縮し（R1.35），不要なKleene星の例と $\mathrm{CCL}_1$ の付録は削除した。
\Location{第9節。}

\Comment{R1.35}
「定義9.1：記号 $\#$ を定義し、第9節を通じて一貫して使用すること。」
「命題9.2は自明で証明不要である。命題9.6も Ginsburg and Greibach, \emph{Deterministic context-free languages}, Information and Control 9 (1966), 620--648 に既に記載されている周知の事実である。」
「補題9.8：Łukasiewicz 言語と Dyck 言語の関係も周知である。」
「命題10.3の証明も省くか大幅に短縮できる。命題10.5の証明では定義3.3との記号の整合を検討されたい。」

\Reply
リセット記号 $\#$ は一度だけ定義して一貫して使う。$D_1$ と Łukasiewicz 言語の関係は Ginsburg--Greibach (1966) および Autebert--Berstel--Boasson (1997) を引用した。投稿時命題10.3の証明は命題9.1に短縮した。決定性文脈自由性はプッシュダウンによる簡潔な議論とし，代入可能性は入口高さ集合 $\mathcal Q(x)$ と脚注の連結恒等式により，$ba$ を含む因子と $a^*b^*$ の因子に分けて直接証明した。他の長い反例証明は命題3.3直後の判定条件または初等的なDFAの記述に縮約した。
\Location{第9節。}

\Comment{R1.36}
「補題9.11は $\sim_h$-代入可能性の基本的性質である。基本的な数学的性質を付録で集めるなら、そこに置いてもよいだろう。」

\Reply
投稿時の基本性質の付録は残さず，正規言語との共通部分・逆準同型像の結果を，消去因子を区別する非空因子の規約のもとで命題3.2にまとめた。残る右商の性質は直後で使う二行の補題9.6で述べ，一事実のための付録は設けない。
\Location{補題9.6。}

\subsection*{投稿時第11節}
\Comment{R1.37}
「$\RS\cap\mathrm{CFL}$ が Clark の合同文脈自由言語（congruential CFL）のクラスに含まれることにも触れる価値がある。」

\Reply
命題9.7として追加した。Clark (2010) は代入可能言語および $(k,\ell)$-代入可能言語を合同文脈自由言語族に含め、Dyck 言語の例も挙げている。改訂稿の命題は、固定 $h$ の各スライスについて包含関係を示し、したがって $\RS(\Gamma)\cap\mathrm{CFL}$ に拡張する。$\{a,b\}$ 上の真の包含は、合同的でありながら $\RS$ の外にある $D_1$ が示す。
\Location{命題9.7。}

% --- END japanese/parts/03_r1_31_37.tex ---


% --- BEGIN japanese/parts/04_r2.tex ---

\section*{Reviewer 2}
学習アルゴリズムの修正については R1.1 と R1.24、以下で問題とされた終端規則の完全性の補修については R2.8 を参照されたい。

\subsection*{固定 $h$ の全般的な位置づけ}
\Comment{R2.1}
「大局的には本論文は、固定長窓を任意の認識可能な合同に置き換えることで Yoshinaka の $(k,\ell)$-代入可能言語の学習を抽象化し、既知の結果を一般化している。この抽象化は強力で十分に理解された枠組みに問題を置く一方、先行研究がクラス全体に一様なパラメータを固定したのに対し、各言語に対して準同型を丸ごと事前に知ることを要求する。準同型の学習が範囲外だと認めているが、$(k,\ell)$-代入可能言語の研究のための抽象化以上に、なぜこの設定が興味深いのか明らかでない。有望な方向は、第8--10節でも触れているプッシュダウン・オートマトンとの関係を深めることだろう。Takada の研究との関係も調べる価値がある。……より実用的な方向では、本研究をオートマトンの型付けを文脈自由代入可能言語へ拡張するものとみなし、型付けが満たすべき条件を調べることもできる。」

\Reply
有限モノイド準同型の指定は $(k,\ell)$ の選択より強い事前バイアスであり，学習器は $h$ 自体を推定しない。ただし一つの $h$ は $\Chcf$ 全体に対して一様に固定される。命題8.4は固定窓のどのクラスにも属さない非正規線形例を示し，系7.5と定理8.3は窓・有限群型付けおよび線形の場合に符号長で多項式の特徴データを保証する。型付けを動かす合併 $\RS$ の学習可能性は主張しない。$h=h_{k,\ell}$ のときの再構成は Yoshinaka (2008, Section 4) の部分語規則方式を特殊化する（表現上の差異は第4.1節）。

序論には Kanazawa (1996; 1998, Chapters 5, 8) の範疇文法の学習を相補的な先行研究として挙げたが，それらから本稿の固定 $h$ 定理が従うとは主張しない。Yoshinaka (2015, Theorem 1) の一般的多ソート枠組みは，固定 $h$ の非空繊維 $h^{-1}(m)\cap\Sigma^+$ をソートとすることで定性的学習可能性を包含する。この型の選択をYoshinaka本人に帰属させる趣旨ではない。有限モノイドに固有の定量的・構造的結果とは区別する。

第3節では Yoshinaka (2008, Proposition 1) の非閉包例と命題3.2を対比し，消去を伴う逆像に対しては型付けを変更して閉包性を得るのであって，元の固定窓や固定 $h$ のクラスの閉包性を主張しないことを明示した。命題3.2(iii)の脚注は消去フラグ $e_\varphi$ の必要性を例示する。局所自明半群の判定条件は命題3.5でPinに帰属させ，補題7.4は局所自明半群と群の直積へのリフトによる窓・有限群の核の特徴づけを与える。系7.5の通常厚みに関する上界は型付き生成語の評価と正規化を利用する。命題6.3・補題7.6・系6.4は，二元半束型付けによる指数的特徴データと型付き厚みの下界を与えるが，任意の型付けについて代数的条件の必要性やGold学習器一般の実行時間・収束時刻下界は主張しない。

序論には Coste, Fredouille, Kermorvant, de la Higuera (2004) の自動機の型付けを引用した。型付けの適切さとは分布上の衝突を有限指数のモノイド合同で分離することであり，構文合同そのものを精密化する条件は正規性を強制する。これは型付けを見つけるアルゴリズムではない。第8節のTakadaの正規制御集合との関係は限定的であり，一般のプッシュダウン特徴づけや同値性は主張しない。
\Location{序論、第3節（命題3.2直前の Yoshinaka との比較）、第4節、命題3.2・3.5、主定理3.4(iv)、命題6.3・系6.4・付録C、第7節・補題7.4・7.6・系7.5・付録A.1、命題8.4、第9.2節（$L_{\mathrm{Yreg}}$ の定義）、第8節の制御集合に関する議論。}

\subsection*{主要な懸念1：外側文脈の型付けと輸送}
\Comment{R2.2}
「最大の懸念は型付き精密化、とくに輸送にある。規則(2)は外側の型に何の制約もないまま、非終端記号を任意の文脈間で移動させる。これは意図されたことなのか。外側の成分 $m,n$ が非終端記号の出現位置を記録するなら、それを移動させる規則が，なぜその成分を完全に無視するのか。私の見るところ、この成分は証明で一度も参照されず、すべての $A^{m,n}_p$ を $A_p$ に置き換えても結果は変わらない。著者は同意するか。そうでなければ、外側の型がどこで参照されるのか示してほしい。これは定義2.3そのものから従うように思える。代数的条件が課されるのは置換される因子だけであり、その周辺文脈ではない。ゆえにその定義を模倣する学習器では、輸送を制限する条件は何もないはずだ。もしそうなら、第1.1節で予告する『型付き輸送』は論文のどこにも登場せず、出現を外側の型に応じて分離すべきだと示す例も存在しえない。自明なモノイドで空虚となる生成語型付けは本当に必要なので、それを脚注のままにしておくべきではない。」

\Reply
同意する。未使用の外側文脈型注釈を削除し，標的側の $A_p^{m,n}$ を生成語型付きの $A_p$ に単純化し，仮説側の $[x:u,v]$ を観測部分語 $x$ ごとの $\NT{x}$ にまとめた。単項規則は標本の同一文脈と $h$-型一致を使い，独立の文脈輸送規則は不要となった。投稿時の外側文脈型の必要性の主張を撤回した。健全性は定理4.2の不変量 $\NT{x}\Rightarrow^*w$ ならば $x\equiv_Lw$ かつ $h(w)=h(x)$ により証明し，外側の型を参照しない。

一方，生成語型付けは本稿の有限証拠集合に基づく完全性証明に必要である。第5.1節に本文で説明し，標準親生成語から型付き終端生成語または子生成語の連結へ進む単項規則(U)を許す。偶奇型付けの例はこの証明での必要性を説明するのであって，任意のアルゴリズムへの普遍的な必要性を主張するものではない。
\Location{第4節、第5.1節、定理5.3。}

\subsection*{主要な懸念2：線形定理の射程}
\Comment{R2.3}
「線形の場合の結果には別の疑問がある。対象のクラスはどれほど大きいのか。末尾の三つの節は $\RS$ の位置づけを述べるだけで、$\mathcal C_h^{\mathrm{lin}}$ について述べておらず、定理7.9が通常の線形代入可能言語の先にどこまで届くのか読者にはわからない。上限付きカウンター言語は正規だから線形でもあり、Yoshinaka の階層との分離はすでに線形クラスの中で成り立つ。一方、$\{a^nb^n\}^*$ は線形ではないため、最後の節は線形のレベルについて何も示していない。現状では、非正規な線形の場合を未解決に残し、射程不明なクラスに対する技術的な強化に見える。」

\Reply
第8節に具体的な非正規線形例と正規フィルターによる構成を加えた。命題8.4の $L_{\pm,e}$ は固定有限モノイド型付けで代入可能な線形・非正規言語だが，Clark--Eyraudクラスにも任意の $(k,\ell)$-クラスにも属さない。命題3.2(ii)と線形言語の正規言語との共通部分の閉包性により，$L\in\Chlin$ かつ $R_{\mathrm{reg}}=g^{-1}(\mathsf{Acc}_g)$ ならば $L\cap R_{\mathrm{reg}}\in\mathcal C_{h\times g}^{\mathrm{lin}}$ である。この例の証明は $L_{\pm,e}=L_{\mathrm{all}}\cap R_{\mathrm{reg}}$ と書き，別に四元モノイドの具体的な型付けも示す。正規言語段階の分離は Yoshinaka (2008, Proposition 1) の既知結果に帰する。

定理8.3は固定 $h$ の代入可能な線形言語に対し，選んだ線形CFGに相対的な集合依存 $\Bh$ の多項式特徴データを保証する。これは一般の線形言語についてのde la Higueraの否定的結果に反しない。$\Ah$ の逐次的な多項式データ保証も主張しない。
\Location{命題3.2(ii)、定理8.3とその適用範囲の補足、命題8.4（証明中の正規フィルターの議論）および後続の四元型付けの段落。}

\subsection*{その他の指摘}
\Comment{R2.4}
「この論文はもっと簡潔にすると改善される。第4節の多くは定義であり、いくつかの結果は直前の証明で示したことを繰り返している。対象を導入する前に使用することも多く、読者は未解決のまま説明を読み進めなければならない。また $\widetilde G$ は $\widetilde G^{\mathrm{full}}$ の刈り込みであるから、記号が到達可能かつ生成可能なのは定義上当然であり、毎回断る必要はなく、その規約を一文で述べれば十分だ。」

\Reply
同意する。到達可能性・生成可能性の繰り返しの限定を削除し，第5.1節で縮約済み精密化の規約を一度だけ述べた。旧補題4.6--4.7，5.3--5.4および定理4.9，4.12などの冗長な結果を削除した（R1.22，R2.9）。再構成作用素は標的側精密化より先に，$\omega$ は初出前に，$\chi$ は証拠集合の直前に，\emph{locally trivial} は命題3.5より先の第2.2節で定義した。定量的議論を第6節（一般の型付き厚み），第7節（窓・有限群），第8節（線形言語）に分け，長い正規化の証明は付録に移した。
\Location{第4節、第5.1節、第6--8節、付録A--B。}

\Comment{R2.5}
「命題3.1：$L=h^{-1}(F)$ という仮定は正規言語を正確に特徴づける。『In particular』ではなく『Thus』とすべきで、証明で理由を明示してもよい。」

\Reply
同意する。命題3.1は有限モノイドによる認識可能性の特徴づけを直接述べ、そこからすべての正規言語が $\RS$ に属すると結論する。標準的な特徴づけについて Eilenberg (1974, Chapter III, Propositions 10.1--10.2) を引用し、必要な帰結も示す。$L=h^{-1}(\mathrm{Acc})$ で $h(x)=h(y)$ なら、すべての文脈 $(u,v)$ に対して $h(uxv)=h(uyv)$ となるから分布は等しい。
\Location{命題3.1。}

\Comment{R2.6}
「定義4.11：基底 $\mathsf B(\widetilde G)$ は使われていない。定理4.12、注4.13とともに導入する必要があるのか。$\mathrm{NT}(\cdot)$ と $\mathrm{Rule}(\cdot)$ も不要であり、$G$ を組として定義し、その $V,P$ を使えば十分だ。」

\Reply
同意する。これらの対象と記法は削除した。
\Location{第5.1節。}

\Comment{R2.7}
「例4.3は SSBNF ではない。」

\Reply
旧例は全面的に削除した。改訂後の証明では不要となったためである。
\Location{第5.1節。}

\Comment{R2.8}
「補題5.2(i)：$X$ に複数の終端規則があるなら $\omega(X)$ だけでは不十分で、規則(3)を使うべきだ。証明を修正する必要がある。」

\Reply
指摘は正しい。投稿時の完全性証明には実際に欠落があった。型付き終端規則 $X\to a$ ごとに同じ標本文脈の証拠語 $u_Xav_X$ を明示的に含め，標準証拠語 $u_X\omega(X)v_X$ と合わせて，必要に応じて単項規則(U)を $\NT{\omega(X)}$ から $\NT{a}$ に適用し，語彙規則(L)を使う。型付き精密化による $h(\omega(X))=h(a)$ がその単項規則を正当化する。
\Location{第5.1節（証拠集合）、定理5.3（終端規則の基本場合）。}

\Comment{R2.9}
「補題5.3と5.4は必要か。」

\Reply
同意する。両補題は削除した。
\Location{第4--5.1節。}

\Comment{R2.10}
「系5.8は『有限モノイド $M$ への固定準同型 $h$ について……』、または『準同型 $h$ が与えられているとき……』という文で始めるべきだ。」

\Reply
改訂後の Gold 同定の文は、提案された『有限モノイド $M$ への固定準同型 $h:\Sigma^*\to M$ について……』という形で始める。
\Location{系5.5。}

\Comment{R2.11}
「第7.1節の標準 SSLNF への変換（3ページ）は、短縮するか付録に移せるだろう。」

\Reply
付録に移す案を採用した。第8節には定理8.3に直接必要な命題8.1・補題8.2を残し，補題8.2の短い証明も本文に置いた。正規化の証明は付録Bにまとめた。
\Location{命題8.1と周辺、補題8.2とその証明、定理8.3とその後の補足、付録B。}

% --- END japanese/parts/04_r2.tex ---

\end{document}
